% WS3 -- Types A+B: electron configurations and PES. Blocks 3-4.
\documentclass{apchem-worksheet}

\apcunit{1}
\apcunittitle{Atomic Structure and Properties}
\apcdoctitle{Worksheet 3 \textbullet\ Configurations \& PES}
\apcsubtitle{CED 1.5, 1.6 \textbullet\ configurations cold, spectra warm}

\begin{document}
\wsheader[Configurations from memory --- no periodic-table peeking until you
check. PES answers must cite peak position or area as evidence.]

\problem[1.5]{Write shorthand ground-state configurations:}
\begin{parts}
  \item \ce{Si}: \answerblank[4cm]{$[\ce{Ne}]\,3s^2\,3p^2$}
  \item \ce{K+}: \answerblank[4cm]{$[\ce{Ar}]$}
  \item \ce{N^3-}: \answerblank[4cm]{$[\ce{Ne}]$ ($1s^22s^22p^6$)}
  \item \ce{Se}: \answerblank[4cm]{$[\ce{Ar}]\,4s^2\,3d^{10}\,4p^4$}
\end{parts}

\problem[1.5]{Classify each as ground state, excited state, or impossible ---
and name the element where one exists:}
\begin{parts}
  \item $1s^2\,2s^2\,2p^2$: \answerblank[4.5cm]{ground --- carbon}
  \item $1s^2\,2s^2\,2p^7$: \answerblank[4.5cm]{impossible (p max 6)}
  \item $1s^2\,2s^2\,2p^5\,3s^1$: \answerblank[4.5cm]{excited neon (10 e$^-$)}
\end{parts}

\problem[1.5]{How many \emph{unpaired} electrons in a ground-state N atom?
Draw the 2p orbital diagram to support your count. \workspace[1.8cm]
\keyonly{\begin{apcanswer}3 --- Hund's rule: $\uparrow\ \uparrow\ \uparrow$
across the three 2p orbitals.\end{apcanswer}}}

\problem[1.6]{The complete PES spectrum of an unknown element:}

\begin{center}
\begin{tikzpicture}
\begin{semilogxaxis}[width=9cm, height=4.4cm, x dir=reverse,
    xlabel={binding energy (MJ/mol)}, ylabel={rel.\ \#\,e$^-$},
    ytick=\empty, xtick={126,9.07,5.31,0.74},
    xticklabels={126, 9.07, 5.31, 0.74},
    ymin=0, ymax=7, axis lines=left]
  \addplot[ycomb, seriesA, line width=2.5pt]
    coordinates {(126,2) (9.07,2) (5.31,6) (0.74,2)};
\end{semilogxaxis}
\end{tikzpicture}
\end{center}

\begin{parts}
  \item Assign every peak to a subshell:
        \answerblank[8.9cm]{126 $\to 1s^2$; 9.07 $\to 2s^2$; 5.31 $\to 2p^6$;
        0.74 $\to 3s^2$}
  \item Identify the element, citing total electron count as evidence:
        \answerblank[6.8cm]{$2+2+6+2 = 12$ e$^-$ $\to$ magnesium}
  \item Why does the 2p peak sit at \emph{lower} binding energy than 2s, even
        though both are in shell 2? \answerlines[2]{2p penetrates less toward
        the nucleus than 2s, so it is shielded slightly more and held
        slightly less tightly.}
  \item Predict two specific changes in this spectrum for the \emph{next}
        element (Al). \answerlines[2]{All peaks shift to higher BE (13
        protons), and a new small peak (area 1, the $3p^1$) appears at lower
        BE than the 3s peak.}
\end{parts}

\problem[1.6]{The 1s binding energy of C is 28.6 MJ/mol; of O it is 54.4
MJ/mol. Explain the difference in one Coulombic sentence --- no octet
language. \answerlines[2]{1s electrons are essentially unshielded, so the
6-proton pull in C vs the 8-proton pull in O acts at similar distance ---
higher $Z$, higher binding energy.}}

\begin{spiralstrip}{Blocks 1--2 \textbullet\ keep it warm}
  \item Average mass, peaks at 85 u (72\%) and 87 u (28\%):
        \answerblank[3cm]{85.6 u (Rb)}
  \item Moles of \ce{Al} in \SI{5.4}{\gram}: \answerblank[3cm]{\SI{0.20}{\mole}}
  \item Empirical formula: 80.0\% C, 20.0\% H:
        \answerblank[3cm]{\ce{CH3}}
  \item If its molar mass is \SI{30}{\gram\per\mole}, molecular formula:
        \answerblank[3cm]{\ce{C2H6}}
  \item One-particle-type diagram with two-atom units of different elements:
        element, compound, or mixture? \answerblank[3cm]{compound}
\end{spiralstrip}

\end{document}
